% ============================================================================
% CHAPTER 2 - LITERATURE REVIEW
% Rubric mapping: CO9 (5 marks) is assessed from 2.1, 2.2, 2.3.
% Works cited here were checked against their sources in September 2026. Where a
% system is a public model or repository rather than a peer-reviewed paper, the
% text says so.
% ============================================================================

\section{Preliminaries}

\subsection{Code-mixing and Banglish}

Code-switching is the use of more than one language by the same speaker in the same
conversation. Two cases are usually distinguished. Inter-sentential switching means the
speaker finishes a sentence in one language and begins the next in another.
Intra-sentential switching, also called code-mixing, means the switch occurs inside a
single sentence. The second case is harder for speech systems, because the utterance
cannot be assigned to one language and handed to one monolingual model
\cite{agro2025cssurvey}.

Bangladeshi classroom speech is strongly intra-sentential. In the recordings used here a
single clause regularly carries a Bengali verb, an English noun phrase and an English
technical term, so the speech is Bengali in grammar and often English in content words.

The written form matters as much as the spoken one. Bengali has its own script, but it is
also widely written in the Roman alphabet in informal digital text, a form usually called
Banglish. There is no spelling standard, so the same word appears as \textit{ami},
\textit{aami} or \textit{amee} depending on the writer. Romanised Bangla is now an object
of study in natural language processing in its own right. BanglaTLit, for example, is a
benchmark of 42,700 romanised Bangla samples with a pre-training corpus of 245,700 more,
built for converting romanised Bangla back into the Bengali script
\cite{fahim2024banglatlit}. Work of this kind establishes romanised Bengali as a real
written form with real users, which is the reason it was chosen here as the output of the
speech model rather than the Bengali script.

\subsection{Speech recognition, error rates and the models used}

Automatic speech recognition takes an audio signal and returns text. Its usual measure is
the word error rate, the number of word substitutions, insertions and deletions needed to
turn the output into the reference divided by the number of reference words. Character
error rate is the same calculation over characters. Both can exceed 100 per cent when the
system produces more text than was spoken.

Two properties of these measures shaped how results are reported in this research. A single
long wrong output for one short clip pushes a mean error rate above 100 per cent and
hides everything else, so medians over clips are reported and systems are compared clip
by clip. Word error rate is also unkind to a language with no spelling standard, because
two acceptable spellings of the same word count as an error; Section \ref{sec:cs-metrics} reviews what
the code-switching literature does about this.

Modern recognisers are built in one of two ways, and the difference matters later in this
research because the two write their output differently. The first is connectionist temporal
classification, which scores an alignment between audio frames and a fixed set of output
symbols \cite{graves2006ctc}. A model of this kind can only ever emit characters that are
in its output alphabet, which is decided before training. The second is the encoder-decoder
transformer \cite{vaswani2017attention}, which generates text one token at a time and can
in principle produce any string its tokeniser can represent. Conformer models add
convolution to the transformer encoder and are a common choice for speech
\cite{gulati2020conformer}.

Most recent recognisers for languages with little transcribed audio start from a model
pretrained on unlabelled speech. wav2vec 2.0 learns from masked audio with a contrastive
objective \cite{baevski2020wav2vec2}, HuBERT predicts hidden units discovered by
clustering \cite{hsu2021hubert}, and WavLM extends the idea to a wider set of speech tasks
\cite{chen2022wavlm}; a review covers the family and how the approaches relate
\cite{mohamed2022sslreview}. The same recipe has been scaled across many languages at once,
first with XLSR \cite{conneau2021xlsr} and then with XLS-R \cite{babu2022xlsr2}, which is
where most published Bengali recognisers begin. WavLM is used directly in this research, not
as a recogniser but to produce a voice embedding for each recording, which is how the
lecturer of every video was confirmed. Where labelled audio is scarce, SpecAugment masks
bands of time and frequency during training and is close to standard practice
\cite{park2019specaugment}.

Whisper is a family of encoder-decoder transformer models trained on about 680,000 hours
of weakly supervised multilingual audio \cite{radford2023whisper}. Its decoder receives
special tokens stating the language and the task, where the task is either transcribe or
translate, which explains two behaviours central to this work. Because the model is told
which single language it is hearing, mixed language audio is pushed towards one of them,
and for Bengali audio containing English words the output is usually English. Because the
decoder is autoregressive with a strong internal language model, it can enter a repetition
loop in which the same phrase is emitted until the token limit is reached. The Whisper
paper proposes a fallback rule for this, in which a hypothesis whose gzip compression
ratio exceeds a threshold is decoded again with sampling \cite{radford2023whisper}, and
WhisperX avoids the problem by segmenting audio with an external voice activity detector
instead of relying on the model's own timestamps \cite{bain2023whisperx}. Loops are
reported in this research as a counted quantity.

A vision-language model takes an image and a text instruction and returns text. Open
models such as Qwen2.5-VL read text inside an image well enough that a separate optical
character recognition stage is often unnecessary \cite{bai2025qwen25vl}. Directing such a
model at a specific part of an image is a separate problem, which Set-of-Mark prompting
solves by drawing marks such as numbered boxes onto the image and asking the model to
answer per mark \cite{yang2023som}. The board reader in this research uses the same idea,
with the boxes found by our own image processing and the numbers carried through into the
generated notes.

An instruction tuned language model such as Qwen2.5-7B-Instruct can be given a transcript
and asked for structured notes \cite{qwen2025qwen25}. The risk is known from the
summarisation literature. Models produce fluent text that the input does not support, and
a study of state of the art abstractive summarisers found that a large share of generated
summaries contain hallucinated content \cite{maynez2020faithfulness}. Faithfulness asks
whether the output follows the input given to the model; factuality asks whether it is
true of the world. A lecture note must be faithful, because a statement that is true in
general but was not part of the lecture is still the wrong note.

\subsection{Fine-tuning, validation, ablation and statistical testing}

Fine-tuning means continuing to train a pretrained model on new data so that it fits a
new domain. Training all weights of a model with hundreds of millions of parameters
requires a large amount of memory, because the optimiser keeps extra state for every
parameter, and it risks destroying what the model already knows when the new dataset is
small.

A family of methods answers this by training a small number of new parameters and leaving
the pretrained weights alone. Adapters insert small trainable layers between the frozen
ones \cite{houlsby2019adapters}, and prefix tuning learns a short sequence of vectors that
is prepended to the input instead of changing any weight \cite{li2021prefix}. Low-Rank
Adaptation is the member of that family used here. A later variant, QLoRA, quantises the
frozen model to four bits so that a much larger model fits on one card
\cite{dettmers2023qlora}; this research did not need it for the recogniser, but the same
idea is what makes the note-writing model practical on a single consumer card.

Low-Rank Adaptation avoids both problems \cite{hu2022lora}. The original weight matrix is
frozen, and beside it two small matrices are trained whose product has the same shape as
the original, so the layer output becomes the sum of the frozen path and the small
trained path. If the original matrix is $d \times k$ and the two small matrices have inner
dimension $r$, the number of trained parameters falls from $d \times k$ to $r(d + k)$, and
$r$ is typically 8, 16 or 32. The frozen model is untouched, so the adapter can be removed
or swapped. In speech work the adapters are usually placed on the attention projections,
most often the query and value projections \cite{yang2025spt}. The rank $r$ controls how
much the adapter can change the model and the learning rate controls the size of each
training step; both are hyperparameters, chosen by the person training the model rather
than learned from the data.

Four further terms recur in Chapters 4 and 5. The training set is what the model learns
from, the validation set is used to choose settings such as the learning rate, and the
test set is used once at the end to report the result. Choosing settings by looking at the
test set makes the reported number optimistic, because the choice has been fitted to data
that is supposed to be unseen. Hyperparameter tuning is the search over settings, and it
is defensible only when scored on validation data, with the order of the search and the
rule for accepting a winner written down in advance. An ablation removes or changes one
component and measures the effect, which is how a claim that a particular part is
responsible for a result can be supported. Finally, in speech work it matters greatly
whether the test speaker was heard during training, so this research reports both an unseen
speaker design and an unseen recording design and labels each one.

Training in this research uses AdamW, which decouples weight decay from the gradient update
and is the usual optimiser for transformer fine-tuning \cite{loshchilov2019adamw}, through
the Hugging Face Transformers library \cite{wolf2020transformers}.

Comparisons here are paired, because both systems are run on the same clips or boards. The
Wilcoxon signed-rank test ranks the differences by size and asks whether positive and
negative differences are balanced \cite{wilcoxon1945}, while the sign test counts only how
many items improved and how many got worse. Both are used, and the direction of the
difference is reported beside every p-value.

\section{Review of Existing Research}

The work this research builds on falls into four areas. The first is speech recognition
for Bengali and for code-switched speech, which supplies both the models and the
evidence that code-switching is where they fail. The second is the small set of systems
that target Bengali-English speech directly, which is where this research has to state
what is new. The third is how code-switched recognition is measured, which matters here
because two of the results in Chapter 5 are about a metric misleading rather than about
a model improving. The fourth is reading a whiteboard and generating lecture notes,
which is the vision half of the pipeline. Each is reviewed in turn, and the section ends
with the gap the four leave open.

\subsection{Bengali and code-switched speech recognition}

Bengali speech recognition is an active area with substantial resources behind it.
OOD-Speech is the largest public Bengali corpus of its kind, with 1177.94 hours of
crowdsourced training audio from 22,645 native speakers and a 23.03 hour test set
annotated by hand from seventeen sources, including television drama, audiobooks, talk
shows and online classes \cite{rakib2023oodspeech}. Its purpose is to measure how badly
models degrade when the test domain differs from the training domain, which is exactly
the failure this research would meet by training on read speech and testing on a classroom.
Two other resources sit behind most published Bengali systems. Common Voice is a
crowdsourced multilingual corpus whose Bengali portion is the training set for a large
share of the public models \cite{ardila2020commonvoice}, and Multilingual LibriSpeech
supplies read speech for the same purpose in other languages \cite{pratap2020mls}. For
Indian languages specifically, IndicSUPERB provides the Kathbath corpus and a benchmark
across twelve languages, which is the closest thing the region has to a standard
evaluation \cite{javed2023indicsuperb}.

On the modelling side, applying wav2vec 2.0 to the Bengali Common Voice dataset gives a
word error rate of 0.2524 on a validation split of 7,747 samples
\cite{shahgir2022wav2vec2bengali}, and several Whisper models fine-tuned on Bengali
Common Voice are publicly available. The Bengali branch of the earlier phase of this
project used one of them, a Whisper-small model fine-tuned on Bengali Common Voice.

All of these write Bengali script. For a classroom whose technical vocabulary is English
that is a problem, not a solution, because the output turns \textit{NAND gate}
into a Bengali spelling of the English sound, which a search for the English term cannot
match and which no longer records the word the lecturer actually used. This is not left as an argument: Section
\ref{sec:baselines} reports what five of these published models actually produce when they
are run on the test clips of this research.

Writing a South Asian language in the Latin alphabet is itself a studied problem. The
Dakshina dataset pairs native-script and romanised text for twelve South Asian languages
and documents how unstandardised the romanised side is \cite{roark2020dakshina}, and
Aksharantar is a large transliteration corpus built to train models that move between the
two \cite{madhani2023aksharantar}. For Bengali specifically, BanglaTLit addresses the
reverse direction, turning romanised Bangla back into Bengali script
\cite{fahim2024banglatlit}. What none of this work does is treat the romanised form as the
transcription target for speech, which is the position taken here. The reason is not that
the romanised form is preferred for reading, but that it is the only target that records a
code-mixed utterance without committing it to one language or the other.

Code-switched recognition has its own literature, surveyed in general terms by Sitaram et
al. \cite{sitaram2019survey} and, for end-to-end systems, in a recent systematic review
that catalogues the languages, datasets, metrics and model choices used across the field
\cite{agro2025cssurvey}. Two observations from that review apply here: intra-sentential
switching is the hard case, and the field is short of natural, spontaneous, in-domain
speech data, because most corpora are read or scripted. The best known exception is SEAME,
about 192 hours of spontaneous Mandarin-English conversation, which has served as the
standard benchmark for the field for over a decade \cite{lyu2010seame}. Nothing of that
kind exists for Bengali-English. Winata et al. also show that a multilingual model is not
automatically a code-switching model, since handling two languages separately is not the
same as handling them in one sentence \cite{winata2021multilingual}.

For South Asian languages, Chadha et al. built code-switched recognition for Indic
languages using wav2vec 2.0 with language identification, reporting word error rates of
21.77 for Hindi-English and 28.27 for Bengali-English \cite{chadha2022indic}. Those
figures are much lower than the ones reported here, but not because their system is
better at this task. Their transcription target is the native script for the Bengali
parts, their audio comes from a different and generally cleaner domain, and their test
split is not made of whole unseen lectures. Error rates across different datasets and
transcription targets are not comparable, and are not presented as such in this research.

On the efficiency side, Yang et al. study parameter efficient adaptation of Whisper for
code-switching, comparing soft prompt tuning variants against LoRA on the SEAME and
ASRU2019 Mandarin-English datasets, and find deep prompt tuning competitive with LoRA
while preserving performance on the original languages \cite{yang2025spt}. Their
conclusion supports the design used here, in which a small adapter on a frozen
multilingual model adds a code-switched variety without retraining everything.

\subsection{The nearest Bengali-English systems}

Four efforts are close enough to this research to be named explicitly, and Table
\ref{tab:neighbours} compares them.

\begin{table}[htbp]
\centering
\caption[Systems that recognise Bengali-English code-switched speech]{Systems that recognise Bengali-English code-switched speech, with the system
built in this research. The error rates come from different datasets with different
transcription targets and are not comparable with one another.}
\label{tab:neighbours}
\small
\begin{tabular}{p{2.8cm}p{2.3cm}p{3.1cm}p{2.0cm}p{2.5cm}}
\toprule
System & Model & Data & Output & Reported error \\
\midrule
MediBeng Whisper Tiny \cite{ghosh2025medibeng} &
Whisper-tiny, full fine-tune &
Synthetic clinical dialogue &
English translation &
WER 0.01, BLEU 0.98 \\
\addlinespace
BEHE-CMDisfl \cite{mitra2026behe} &
Kaldi GMM-HMM baseline &
Synthetic speech from LLM text and TTS, 1.3 h subset &
Not stated &
WER 37.74\% \\
\addlinespace
Indic code-switched ASR \cite{chadha2022indic} &
wav2vec 2.0 with language ID &
Indic code-switched corpora &
Native script &
WER 28.27 (Bengali-English) \\
\addlinespace
Whisper-small-\newline benglish (public
model, no paper) &
Whisper-small fine-tune &
194 clips of about 30 s, roughly 1.6 h &
Not documented &
WER 0.4123 on a random 10 per cent split \\
\addlinespace
This research &
Whisper-large-v3-turbo with LoRA &
5.15 h of hand-checked classroom speech, 28 lectures, 3 lecturers &
Romanised Banglish &
CER 67.7 to 15.8, WER 93.9 to 41.7, six whole unseen lectures \\
\bottomrule
\end{tabular}
\end{table}

MediBeng is a translation system, not a transcription system, and its training
data is synthetic \cite{ghosh2025medibeng}; a word error rate of 0.01 with a BLEU score
of 0.98 is what synthetic data produces and should not be read as accuracy on real
clinical speech. BEHE-CMDisfl is likewise synthetic, generated by prompting a language
model for code-mixed text with disfluencies and rendering it with text-to-speech
\cite{mitra2026behe}; it is a useful resource and makes the same point about missing
spontaneous data, but it is not spontaneous speech. The Indic work uses real speech but
targets the native script \cite{chadha2022indic}. The community model
Whisper-small-benglish is closest in intent, since its card names lecture notes
among its uses, but it is a model card and not a study. It is trained on about 1.6
hours and evaluated on a random ten per cent of the same clips, so the same speakers and
very likely the same recordings appear on both sides of the split, and no example output
is given, so its orthography cannot be confirmed. A public repository by Asrarul applies
LoRA to Whisper-large-v3 for Bengali-English meeting speech and is the closest match in
method, but reports no results.

The conclusion is narrower than a claim of the first Banglish speech recogniser.
Bengali-English code-switched recognition exists. What could not be found is a system
that transcribes spontaneous Bengali-English classroom speech into romanised Banglish,
evaluated on whole lectures held out from training against hand-checked references, with
the whiteboard read alongside it. Every claim of novelty in this research is written in
that form.

\subsection{What the published models produce on our data}
\label{sec:baselines}

The error rates in Table \ref{tab:neighbours} come from different datasets with different
transcription targets, so they cannot be compared with one another or with this work. An
argument of that kind is weaker than a measurement, so five published models were run on
the 177 held-out clips of this research, with the same references, the same normaliser and
the same metric code that produced the result in Chapter 5. Table \ref{tab:baselines}
gives the outcome.

Two decisions keep the comparison fair. Each model decodes under its own generation
settings, with no language or task token imposed on it, because a Bengali model was
trained to emit Bengali and forcing it into another mode would break it for reasons that
would be our fault. And because our references are romanised, a model that answers in
Bengali script scores badly no matter how well it heard the speech, so the alphabet each
model wrote in is counted and reported beside its error rate.

\begin{table}[htbp]
\centering
\caption[Published Bengali models run on the 177 held-out clips]{Published Bengali and Bengali-English models run on the 177 held-out clips of
this research. Median per-clip rates. The transliterated column romanises each model's
Bengali output before scoring and keeps the better of two schemes per clip, so it is an
upper bound on what the model could score if orthography were free.}
\label{tab:baselines}
\small
\begin{tabular}{p{4.3cm}p{1.5cm}p{1.9cm}p{1.4cm}p{3.3cm}}
\toprule
Model & CER & CER after translit. & WER & What it wrote \\
\midrule
tugstugi Bengali Whisper-medium & 100.0 & 47.1 & 100.0 & Bengali script, all 177 clips \\
\addlinespace
Whisper-small-benglish & 89.9 & 66.0 & 95.8 & Bengali script 124, mixed 53 \\
\addlinespace
bangla-ASR-v5 & 98.6 & 73.0 & 100.0 & Bengali script, all 177 clips \\
\addlinespace
BanglaASR & 98.6 & 73.8 & 100.0 & Bengali script, all 177 clips \\
\addlinespace
MediBeng Whisper-tiny & 109.2 & n/a & 146.5 & English translation, 70 runaways \\
\addlinespace
Whisper-large-v3-turbo, off the shelf & 67.7 & n/a & 93.9 & English translation \\
\addlinespace
This research & \textbf{15.8} & n/a & \textbf{42.5} & Romanised Banglish \\
\bottomrule
\end{tabular}
\end{table}

Not one of the five produced romanised Banglish on a single clip. Four write Bengali
script and the fifth translates into English. The clearest case is the model whose card
advertises Benglish, which does switch between languages but writes the Bengali half in
Bengali letters: where the lecturer said \textit{hello everyone, welcome to the second
class of digital logic design}, it returned \textit{Hello} followed by the English word
\textit{everyone} spelled out in Bengali script, then the English clause, then Bengali
script again. The strongest Bengali model in the table writes Bengali script so
consistently that none of its output survives conversion to plain ASCII at all, which is
why both of its error rates are exactly 100.

The obvious objection is that this punishes a model for its alphabet rather than for its
hearing, so the Bengali output was transliterated to Roman and scored again. That closes
much of the gap and it must be reported honestly. The strongest model improves from 100.0
to 47.1, which is better than off-the-shelf Whisper, so it would be wrong to say these
models are no better than an untuned system. It does not close the comparison: 47.1
against 15.8 is still three times the error, and word error rate barely moves at all,
from 100.0 to 90.2, because transliteration can fix an alphabet but cannot invent the
spelling of a language that has no spelling standard.

The conclusion is therefore narrow and measured. These are capable models being asked for
something they were not built to produce. No published system was found that outputs
romanised Banglish, and the transcripts they do produce cannot be read by the student this
research is written for. One candidate, a wav2vec 2.0 Bengali model, was not run; its output
alphabet was inspected instead and is dominated by Bengali characters, but it is not
claimed here as measured.

\subsection{Measuring code-switched recognition}
\label{sec:cs-metrics}

Evaluation is a recognised research topic in this area, and it bears directly on two
results reported in Chapter 5. Hamed et al. benchmarked evaluation metrics for
code-switching recognition against human judgements and found that transliteration
followed by text normalisation agrees most closely with human acceptance
\cite{hamed2022benchmarking}. Emond et al. proposed a transliteration based word error
rate for the same reason, mapping mixed script output into a single script before scoring
\cite{emond2018transliteration}. PolyWER accepts several valid forms of a reference,
including transliterations and translations, on the grounds that plain word error rate is
too strict for speech that can be correctly written in more than one way
\cite{kadaoui2024polywer}.

The common thread is that a measure built around one script or one vocabulary misjudges a
system that legitimately produces another. This research reaches the same conclusion from a
different direction. Its inherited term based measure counts a fixed list of English
words, so a model that transcribes the mixed speech faithfully emits fewer of them and
scores worse for being right. Section 5.4.2 quantifies the effect.

\subsection{Whiteboard content and lecture note generation}

Extracting what is written on a board from a fixed camera lecture video is an established
line of work in document analysis. Yeh et al. extract handwriting and produce a video
summary by statistical modelling and clustering of board pixels, robust to noise,
lighting change and occlusion \cite{yeh2016robust}. Davila and Zanibbi build a
spatio-temporal index of content regions and segment the video by detecting conflicts
between them, producing static image summaries of a whole lecture from a minimal set of
frames with 96.28 per cent recall of connected components \cite{davila2017whiteboard}.
Kota et al. generalise this with a deep network that detects handwritten text, formulae
and sketches and then binarises them, reaching an f-measure of 94.32 per cent for binary
connected components on the AccessMath dataset \cite{kota2019generalized}.

Reading text out of an image has itself changed since that work was published. Document
understanding first moved to models that combine text, layout and image features, such as
LayoutLM \cite{xu2020layoutlm}, and then to models that skip optical character recognition
altogether and map the image straight to the output string, such as Donut
\cite{kim2022donut}. General vision-language models followed. Visual instruction tuning
produced open models that answer questions about an image in natural language
\cite{liu2023llava}, closed models of the same kind read text in photographs well
\cite{openai2023gpt4}, and a recent survey covers how these models are now applied to
visually rich documents and where they still fail \cite{ding2025vrdusurvey}. A whiteboard
is a harder case than a scanned page, because the writing is handwritten, unevenly lit,
partly occluded by the person writing it, and never captured straight on.

These systems solve a closely related problem and solve it well. Two differences matter
here. Their output is a binarised summary image or an index of handwriting regions, which
is the right target for search and navigation, whereas the target in this research is a
photograph-like board that a student can look at inside a note, with the original strokes
and colours intact. None of them reads the content into text with a model that
understands what it is reading, because that capability did not exist when most of them
were published. The board method used here is a tiled mosaic over erase separated eras,
which is engineering built from known parts and not a new algorithm, and Chapter 4
describes it as such.

Turning a recorded lecture into notes is an old goal given new life by large models. M3AV
is a recent academic lecture dataset of nearly 367 hours across five sources with human
annotation of slide text and spoken words, built to support contextual speech
recognition, slide generation and script generation \cite{chen2024m3av}. NoteIt converts
instructional videos into interactable notes by extracting hierarchical structure and
multimodal key information, and evaluates the result both on technical metrics and with a
user study of 36 participants \cite{zhao2025noteit}. Two gaps appear when this work is
placed beside the classroom targeted here. Existing systems assume slides, which are text
that can be read reliably or taken directly from a file, whereas our lecturers write on a
whiteboard and stand in front of what they write. The datasets and systems are also built
for English or Chinese and for speech that stays in one language. The user study in
NoteIt points at something this research has not done. Whether notes are good is a question
about readers, and it is answered by asking readers.

Results from vision-language models in this area are also strongly affected by the
prompt. Promptception, a systematic study of prompt sensitivity across ten large
multimodal models, reports accuracy swings of up to 15 per cent from minor changes in
phrasing and structure \cite{ismithdeen2025promptception}. That finding is the published
context for one of the results in Chapter 5, in which changing the board reading prompt
was worth far more than changing the size of the model.

\section{Summary of Key Findings}

The review supports four conclusions. The first is that the speech problem is real and is
not solved by anything available. Bengali speech
recognition is mature but writes a script that discards the code-mixing a faithful
transcript has to keep. Code-switched
recognition for Bengali and English exists, but the published systems either translate
into English, train on synthetic speech, or target the native script, and the one
community model that matches the intent most closely is undocumented and evaluated on a
random split of its own clips.

The second is that the board problem is well studied, and this research should not
claim otherwise. Extracting
handwritten content from fixed camera lecture video has a long literature with strong
results on public datasets. What is new here is not the extraction but what follows it,
which is reading the rebuilt board with a vision-language model and carrying numbered
regions of it into a generated note.

The third is that measurement is an open problem in this area and not a detail. Several independent
groups have found that standard word error rate misjudges code-switched output and have
proposed transliteration based or multi-reference alternatives. The same class of problem
appears twice in this research, once in the inherited term measure and once in the note
quality measure.

The last is that lecture note generation with large models is active but assumes a
different classroom.
Existing systems expect slides and a single language, and nothing found combines
code-mixed speech, a handwritten whiteboard and generated notes in two languages. Table
\ref{tab:gap} states the gap in the form the rest of the research uses.

\begin{table}[htbp]
\centering
\caption{The gap this research addresses.}
\label{tab:gap}
\small
\begin{tabular}{p{2.8cm}p{5.4cm}p{5.4cm}}
\toprule
Component & What already exists & What this research adds \\
\midrule
Speech &
Bengali ASR in native script; Bengali-English code-switched systems on synthetic or out
of domain data &
Romanised Banglish transcription of spontaneous classroom speech, tested on whole unseen
lectures against hand-checked references \\
\addlinespace
Evaluation &
Transliteration based and multi-reference metrics for code-switched ASR &
A measured demonstration, in two places, that English-string measures reward a
translating model and penalise a transcribing one \\
\addlinespace
Whiteboard &
Handwriting extraction, binarisation and summary images from fixed camera video &
A reconstruction in real camera pixels, read as text by a vision-language model and
scored against hand-verified records \\
\addlinespace
Notes &
Slide based lecture summarisation and note generation in English &
Notes in two languages that point at numbered regions of the board, with quotations
checked word for word against the transcript \\
\bottomrule
\end{tabular}
\end{table}
